\section{Introduction}
\label{sec:1}

A logistics disruption, whether a supplier failure, a warehouse closure, or a demand surge, often has effects that last several weeks or even months. Its cost depends on the severity of the disruption and, just as much, on how quickly and how appropriately operations managers react \citep{sheffi2005book,ivanov2019ijpr}. A robust network absorbs the disruption without degrading. When this is not possible, its resilience, that is, its ability to limit the loss and then recover, depends largely on the decisions taken during the crisis.

Mitigation and contingency strategies are well documented \citep{tomlin2006,chopra2004}. Resilience is classically represented by a loss and recovery curve \citep{bruneau2003,rahiel2026in4pl}. It remains difficult, however, to answer two questions that operations managers ask in practice. Faced with a given disruption, which reaction is most likely to reach the target resilience level? After a poor recovery, was the loss due to the disruption, the network, or the reaction? Answering them requires observations that link a disruption, the decisions taken, and the resulting evolution of the service level. Company data are confidential and incomplete, public databases describe disruptions without the decisions, and simulation studies report only a small number of cases \citep{ramzy2022,do2024}. A deeper difficulty adds to this. A real crisis is experienced only once, under a single management, and what another reaction would have produced remains unknown. Because operations managers react more strongly when the crisis is more severe, the observed effect of a decision is mixed with that of the disruption.

This work removes this difficulty through counterfactual replay: each simulated crisis is replayed identically, first with no reaction and then under several management policies, so that any difference in trajectory is due to the decision alone. Moreover, no existing database offers this property. It is this property that allows a causal model to recommend a decision, and not merely to observe that a decision and an outcome often go together. To overcome this second difficulty, the open digital twin ISOMORPH \citep{zhang2026isomorph} was taken up and extended in a platform developed for this work, named Isomorph-Reborn.

The approach consists of three sequential blocks, as shown in Fig.~\ref{fig:d1}. The first, Simulate, produces labeled crises, that is, crises whose cause (type, target, date, duration, and severity of the disruption) and the decisions taken in response are known exactly. In this block, the multi-echelon logistics network simulated by Isomorph-Reborn is subjected to controlled disruptions to which a simulated operations manager reacts. The second, Measure, summarizes each daily service level trajectory by a seven-parameter resilience curve, from which six ``network resilience'' indicators are derived. The third, Learn, builds a causal Bayesian network on the resulting dataset. This network is used for prognosis, diagnosis, and recommendation, based on the counterfactual concept developed by the 2011 Turing Award laureate Judea Pearl \citep{pearl2009}.

\begin{center}
\resizebox{0.88\textwidth}{!}{% Figure 1: three-block approach (simplified version).
% The former detailed figure remains available: figures/Fig1_approach_EN.pdf

% --- Adjustable block dimensions -----------------------------------------
\def\bw{5.0cm}   % block width
\def\bh{3.0cm}   % block height
% -------------------------------------------------------------------------

\begin{tikzpicture}[
  font=\small,
  bloc/.style={draw=#1, line width=1pt, rounded corners=4pt, fill=#1!6,
               minimum width=\bw, minimum height=\bh},
  titre/.style={font=\small\bfseries, text=#1!80!black},
  etape/.style={draw=black!25, rounded corners=2pt, fill=white,
                text width=3.6cm, align=center, minimum height=0.62cm,
                inner sep=2pt, font=\scriptsize},
  fl/.style={-{Stealth[length=2.4mm]}, line width=1.1pt, black!55},
  pt/.style={-{Stealth[length=1.6mm]}, black!35}]

  % ---- Block 1 ---------------------------------------------------------
  \node[bloc=polreac] (B1) at (0,0) {};
  \node[titre=polreac, rotate=90, anchor=north]
        at ([xshift=7pt]B1.west) {1. Simulate};
  \node[etape] (a1) at (0.40, 0.85) {Network and labeled\\disruption};
  \node[etape] (a2) at (0.40, 0.00) {Crisis replayed: no reaction\\and 3 management profiles};
  \node[etape] (a3) at (0.40,-0.85) {Service level $\mathrm{PI}(t)$ over time};
  \draw[pt] (a1) -- (a2); \draw[pt] (a2) -- (a3);

  % ---- Block 2 ---------------------------------------------------------
  \node[bloc=palmec] (B2) at (5.6,0) {};
  \node[titre=palmec, rotate=90, anchor=north]
        at ([xshift=7pt]B2.west) {2. Measure};
  \node[etape] (b1) at (5.95, 0.45) {Associated 7-parameter resilience curves};
  \node[etape] (b2) at (5.95,-0.45) {6 network resilience\\indicators};
  \draw[pt] (b1) -- (b2);

  % ---- Block 3 ---------------------------------------------------------
  \node[bloc=polprud] (B3) at (11.2,0) {};
  \node[titre=polprud, rotate=90, anchor=north]
        at ([xshift=7pt]B3.west) {3. Learn};
  \node[etape] (c1) at (11.55, 0.45) {Causal Bayesian network};
  \node[etape] (c2) at (11.55,-0.45) {Prognosis, diagnosis, and recommendation};
  \draw[pt] (c1) -- (c2);

  % ---- Sequence --------------------------------------------------------
  \draw[fl] (B1.east) -- (B2.west);
  \draw[fl] (B2.east) -- (B3.west);
\end{tikzpicture}
}
\captionof{figure}{Three-block approach: simulate crises under four policies, measure their resilience, and learn a causal Bayesian network.}\label{fig:d1}
\end{center}

Concretely, this is a parametric platform for generation, analysis, and decision support. The generator of labeled and counterfactually replayed crises applies to any multi-echelon network. It makes inventory cover a factor of study rather than an obstacle. This generator offers a characterization of resilience whose parameters read directly in logistics terms. Finally, it builds a causal model whose structure follows the temporal order of a crisis and in which the management policy is an intervention, which makes its effect on resilience identifiable.

The next section positions this work in the literature. Section~3 describes the approach and Section~4 presents its results, which are discussed in Section~5 before the conclusion. The technical developments are gathered in a supplementary material, in particular to allow genuine repeatability by anyone. Its sections, figures, and tables carry the prefix S (Section~S3, Figure~S1, Table~S1). A reference of this form therefore always points to this supplementary material.
